% ======================= 10. EVALUATION AND CONCLUSIONS =======================
\section{Evaluation and Conclusions}\label{sec:eval}

\subsection{Strengths}

\textbf{One model, four answers.} The gravity model of Q1 is not a decorative first
section; it is the demand engine for the rebalancing program, the incentive supply
curve and the parking forecast. Because the same estimated object is used four times,
the answers are comparable: the 3.48 CNY that a truck costs per bicycle and the 1.20
CNY that a rider is paid are prices in the same market, and the recommendation that
follows --- incentives for the cheap half, trucks for the rest --- is a statement about
where those two prices cross.

\textbf{Identification is treated as part of the modelling.} Three quantities that the
data do not contain are handled explicitly: the sampling rate is estimated and shown to
be irrelevant to the fleet size; the initial inventory is bracketed between two
estimators and the spread is reported rather than hidden; the incentive parameters are
declared unidentified and swept. A reader can therefore see exactly which conclusions
rest on data and which rest on assumption.

\textbf{Validation is both internal and external.} The behavioural model is tested on
two held-out days against two nested benchmarks. The spatial structure is tested
against an independent source --- 1,097 subway entrances from OpenStreetMap --- and the
attractors the model finds are twice as likely to sit near one as a random cell.

\subsection{Weaknesses}

\textbf{Attraction is endogenous.} $A_j$ is measured from the arrivals it already
receives, so the model describes an equilibrium rather than a structural choice model.
It can answer ``what if the distance decay changed?'' but not ``what if a new mall
opened?''. A structural alternative --- a random-utility model with employment and
retail density as attributes --- would answer the second question but requires land-use
data we do not have.

\textbf{The incentive parameters are assumed.} No reward was offered in May 2017, so
$\kappa,\rho,c_0$ cannot be estimated from this dataset. We have been explicit about
this and have shown that the qualitative conclusion survives a factor-of-two error in
each. But the recommended reward of 1.20 CNY is a prediction of the model, not a
measurement, and it should be treated as the design of a pilot rather than a settled
price.

\textbf{The truck fleet is continuous.} Vehicle routing with integer trucks, time
windows and driver shifts is a harder problem than the one solved here. At the scale of
60 vehicles the rounding error is small relative to the uncertainty in the cost
parameters, but an operator scheduling individual trucks would need the integer
program.

\textbf{One reward for everyone.} A personalised reward could extract more surplus and
would shift the supply curve, but it raises fairness and gaming problems that a uniform
reward avoids. The uniform design is a deliberate choice, not a limitation of the
method.

\textbf{No weather, no holidays, no long run.} The seven typical days are mild spring
weekdays. Rain reduces cycling sharply and would change both the rates and the decay,
and the model has nothing to say about it.

\subsection{What would change the conclusions}

The recommendation to rely on incentives for the first half of the deficit is driven by
the ratio of the reward to the truck cost. If riders valued a yuan at less than about
half of what we assume, the optimal reward would fall below 0.6 CNY and the incentive
would cover less than a third of the deficit; if they valued it twice as much, the
incentive would cover two-thirds. The truck-based conclusion of Q2 is more robust: a
four-fold error in the truck price moves the relocation volume by 15\%.

The oversupply diagnosis is the one conclusion that would be overturned wholesale if
the initial inventory were very different from what we estimate. Our defence is that
the zone-level stocks are large compared with the daily fluctuations (the worst
cumulative excursion is 349 bicycles against a minimum zone stock of 1,986), so the
diagnosis survives any plausible redistribution of midnight bicycles among zones.

\subsection{Conclusions}

Beijing's dockless bicycle system in May 2017 was not failing for want of bicycles. It
was failing because the bicycles were in the wrong places, because nobody was paying
for them to be moved, and because the space they occupied was free to the firms that
put them there. The models in this paper make each of those statements quantitative.

The behavioural model shows that demand is a superposition of five behaviours,
dominated by short multi-purpose riding, and that distance explains most of where
riders go with a decay that varies by a quarter over the day. The rebalancing program
shows that the resulting imbalance can be repaired by moving 24,000 bicycles overnight
--- 1.1\% of daily trips --- and that the fleet itself is about a quarter larger than the
service requires. The incentive analysis shows that the first half of the repair can be
bought from riders for 1.2 CNY each, a third of what a truck costs, and that the second
half cannot be bought at any price because the riders are not there. The fence analysis
shows that 626 fences would capture three quarters of arrivals but that the capacity
they need is so concentrated that uncapped fences would simply formalise the pile.

The unifying finding is that the instruments are complements with different domains:
incentives where the traffic is dense, trucks where it is sparse, fences where the
parking must be bounded, and data --- shared, city-wide data --- because none of these
instruments can be aimed without it.

\emph{Every number quoted in this paper is written to a result file by the script that
computes it; the whole chain is reproducible from the raw order file.}
